% =====================================================================
\section{Exercise 2 -- DCC dynamics on cryptocurrency returns}\label{sec:ex2}
% =====================================================================

We work with the daily log-returns of three crypto assets
(BTC, ETH, ADA) loaded from \file{crypto.xlsx}. The sample spans
$T=1827$ trading days from 2 January 2020 to 1 January 2025. The whole pipeline
is in \file{exercise2/} (orchestrated by \file{exercise2/main.py}).

% ---------------------------------------------------------------
\subsection{2.1 -- Univariate descriptives and sample correlations}
% ---------------------------------------------------------------

\begin{table}[H]\centering\small
\caption{Daily log-return descriptives. Skewness uses $b_1$,
kurtosis is the \emph{regular} (non-excess) version.}\label{tab:21-desc}
\begin{tabular}{lrrrr}
\toprule
Asset & Mean & Std & Skew & Kurt\\\midrule
BTC & $0.001404$ & $0.034701$ & $-1.6580$ & $29.3539$\\
ETH & $0.001780$ & $0.044965$ & $-1.3830$ & $23.2127$\\
ADA & $0.001777$ & $0.052556$ & $-0.2361$ & $11.7605$\\
\bottomrule
\end{tabular}
\end{table}

\paragraph{Sample correlation matrix.}
$\rho_{\textsc{btc-eth}}=0.8355$,\quad
$\rho_{\textsc{btc-ada}}=0.6928$,\quad
$\rho_{\textsc{eth-ada}}=0.7353$.

The three series exhibit crypto stylised facts: near-zero mean,
markedly negative skew (especially BTC and ETH), and a pronounced excess
kurtosis (Pearson kurtosis $\gg 3$).

% ---------------------------------------------------------------
\subsection{Theoretical framework: GARCH/GJR + DCC}
% ---------------------------------------------------------------

\paragraph{Marginal volatility -- GARCH(1,1) (slides eq.\,67).}
For a zero-mean return $r_t$,
\begin{equation}\label{eq:garch}
r_t=\sqrt{h_t}\,z_t,\qquad
h_t=\omega+\alpha r_{t-1}^{2}+\beta\,h_{t-1},\qquad z_t\sim\mathcal{N}(0,1)\ \text{i.i.d.}
\end{equation}
Conditions: $\omega>0$, $\alpha,\beta\ge 0$, persistence $\alpha+\beta<1$. The
Gaussian (quasi-) MLE log-likelihood is
\[
\ell(\theta)=-\tfrac12\sum_{t=1}^T\Bigl[\ln(2\pi)+\ln h_t+\frac{r_t^2}{h_t}\Bigr],
\qquad h_1\equiv\widehat{\sigma}_r^{\,2}.
\]
We optimise by L-BFGS-B in \code{exercise2.garch.estimate\_garch}.

\paragraph{Marginal volatility with leverage -- GJR(1,1) (slides eq.\,71).}
\begin{equation}\label{eq:gjr}
h_t=\omega+\bigl(\alpha+\gamma\,\indic\{r_{t-1}<0\}\bigr)r_{t-1}^{2}+\beta\,h_{t-1},
\end{equation}
with persistence $\alpha+\beta+\gamma/2<1$ (because $\Pr(r<0)\approx 1/2$ under
symmetric innovations). $\gamma>0$ captures the leverage effect.

\paragraph{From margins to standardised innovations.}
Let $\hat\eta_{it}$ denote the standardised residual of asset $i$ at time $t$:
\[
\hat\eta_{it}=
\begin{cases}
\Phi^{-1}\!\bigl(\hat F_{i}^{\text{EDF}}(r_{it})\bigr)& \text{(EDF margin, 2.2a)}\\[2pt]
r_{it}/\sqrt{\hat h_{it}}, & \text{(GARCH/GJR margin, 2.2b--c)}.
\end{cases}
\]

\paragraph{DCC dynamics (Engle, 2002 -- slides eqs.\,68--69).}
With $z_t=\hat\eta_t$ (an $n\times 1$ vector of standardised innovations) we
posit
\begin{align}
Q_t &=(1-a-b)\,\hat\Omega+a\,z_{t-1}z_{t-1}^\top+b\,Q_{t-1},\quad Q_1=\hat\Omega,\label{eq:dcc-Q}\\
\Psi_t &=\Lambda_t^{-1/2}\,Q_t\,\Lambda_t^{-1/2},\qquad \Lambda_t=\diag(Q_t),\label{eq:dcc-psi}
\end{align}
where $\hat\Omega=\frac{1}{T}\sum_t z_tz_t^\top$ and $a,b>0$, $a+b<1$. The
normalisation in \eqref{eq:dcc-psi} ensures $\Psi_t$ is a correlation matrix.

\paragraph{Gaussian DCC log-likelihood (slides eq.\,72).}
The copula likelihood (margins concentrated out) is
\begin{equation}\label{eq:dcc-LL-G}
\ell^G(a,b)=-\tfrac12\sum_{t=1}^{T}\Bigl[\ln|\Psi_t|+\eta_t^\top(\Psi_t^{-1}-I_n)\eta_t\Bigr].
\end{equation}

\paragraph{Student-$t$ DCC log-likelihood (slides eq.\,76).}
With $A(\nu)=\dfrac{\Gamma\!\bigl(\tfrac{\nu+n}{2}\bigr)\,\Gamma\!\bigl(\tfrac{\nu}{2}\bigr)^{n-1}}
{\Gamma\!\bigl(\tfrac{\nu+1}{2}\bigr)^{n}}$,
\begin{align}
\ell^t(a,b,\nu)
&=T\ln A(\nu)-\tfrac12\sum_{t}\ln|\Psi_t|
   -\tfrac{\nu+n}{2}\sum_t\ln\!\Bigl(1+\tfrac{1}{\nu}\eta_t^\top\Psi_t^{-1}\eta_t\Bigr)\notag\\
&\quad+\tfrac{\nu+1}{2}\sum_t\sum_{i=1}^n\ln\!\Bigl(1+\tfrac{\eta_{it}^2}{\nu}\Bigr),
\label{eq:dcc-LL-t}
\end{align}
with $\eta_{it}=t_\nu^{-1}(\hat u_{it})$ obtained from the EDF margins. Both
likelihoods are maximised by Nelder--Mead in \code{exercise2.dcc}.

% ---------------------------------------------------------------
\subsection{2.2 -- DCC under three marginal specifications}
% ---------------------------------------------------------------

\paragraph{(a) EDF margins.} $\hat a=0.012308$, $\hat b=0.987592$, $\hat a+\hat
b=0.999900$, $\ln L=1954.19$, with unconditional correlations
$\hat\Omega_{\textsc{btc-eth}}=0.8248$, $\hat\Omega_{\textsc{btc-ada}}=0.7039$,
$\hat\Omega_{\textsc{eth-ada}}=0.7418$. The persistence is essentially 1.
See Figure~\ref{fig:22a}.

\begin{figure}[H]\centering
\includegraphics[width=0.85\linewidth]{fig_2_2a_dcc_edf.png}
\caption{Dynamic correlations under DCC-Gaussian with EDF margins.}\label{fig:22a}
\end{figure}

\newpage
\paragraph{(b) GARCH(1,1)-N margins.} The marginal estimates are reported in
Table~\ref{tab:22b-margins}. The DCC stage gives $\hat a=0.060304$,
$\hat b=0.900190$, $\hat a+\hat b=0.960494$, $\ln L=1964.18$.

\begin{table}[H]\centering\small
\caption{GARCH(1,1)-N marginal estimates.}\label{tab:22b-margins}
\begin{tabular}{lrrrrr}
\toprule
Asset & $\omega$ & $\alpha$ & $\beta$ & $\alpha+\beta$ & $\ln L$\\\midrule
BTC & $4.88\times 10^{-5}$ & $0.1146$ & $0.8595$ & $0.9741$ & $5324.93$\\
ETH & $3.95\times 10^{-5}$ & $0.1030$ & $0.8850$ & $0.9880$ & $4943.94$\\
ADA & $1.11\times 10^{-4}$ & $0.1365$ & $0.8336$ & $0.9701$ & $4633.55$\\
\bottomrule
\end{tabular}
\end{table}

\begin{figure}[H]\centering
\includegraphics[width=0.85\linewidth]{fig_2_2b_dcc_garch.png}
\caption{DCC dynamic correlations with GARCH(1,1)-N margins.}\label{fig:22b}
\end{figure}

\paragraph{(c) GJR(1,1)-N margins.}
Table~\ref{tab:22c-margins}. DCC parameters: $\hat a=0.062861$, $\hat b=0.902904$,
$\hat a+\hat b=0.965765$, $\ln L=1934.41$.

\begin{table}[H]\centering\small
\caption{GJR(1,1)-N marginal estimates. $\gamma$ is the leverage coefficient.}
\label{tab:22c-margins}
\begin{tabular}{lrrrrr}
\toprule
Asset & $\alpha$ & $\beta$ & $\gamma$ & Persistence & $\ln L$\\\midrule
BTC & $0.0713$ & $0.8426$ & $0.0996$ & $0.9637$ & $5340.71$\\
ETH & $0.0852$ & $0.8776$ & $0.0457$ & $0.9857$ & $4946.55$\\
ADA & $0.1288$ & $0.8280$ & $0.0265$ & $0.9701$ & $4634.15$\\
\bottomrule
\end{tabular}
\end{table}

\begin{figure}[H]\centering
\includegraphics[width=0.85\linewidth]{fig_2_2c_dcc_gjr.png}
\caption{DCC dynamic correlations with GJR(1,1)-N margins.}\label{fig:22c}
\end{figure}

% ---------------------------------------------------------------
\subsection{2.3 -- Comparison of dynamic correlation paths}
% ---------------------------------------------------------------

Figure~\ref{fig:23} overlays the three correlation paths for the
three marginal specifications. Table~\ref{tab:23-summary} summarises the
sample distribution of each path.

\begin{table}[H]\centering\small
\caption{Distribution of estimated dynamic correlations across $T=1827$ days.}
\label{tab:23-summary}
\begin{tabular}{llrrrr}
\toprule
Pair & Margin & Mean & Std & Min & Max\\\midrule
\multirow{3}{*}{BTC--ETH}
 & EDF      & $0.8333$ & $0.0538$ & $0.6857$ & $0.9248$\\
 & GARCH-N  & $0.8200$ & $0.0700$ & $0.5129$ & $0.9746$\\
 & GJR-N    & $0.8148$ & $0.0751$ & $0.4887$ & $0.9731$\\\midrule
\multirow{3}{*}{BTC--ADA}
 & EDF      & $0.7187$ & $0.0917$ & $0.4114$ & $0.8697$\\
 & GARCH-N  & $0.7012$ & $0.1051$ & $0.1564$ & $0.9751$\\
 & GJR-N    & $0.6962$ & $0.1132$ & $0.1207$ & $0.9761$\\\midrule
\multirow{3}{*}{ETH--ADA}
 & EDF      & $0.7511$ & $0.0779$ & $0.4928$ & $0.8839$\\
 & GARCH-N  & $0.7410$ & $0.0977$ & $0.3218$ & $0.9706$\\
 & GJR-N    & $0.7388$ & $0.1042$ & $0.2796$ & $0.9692$\\
\bottomrule
\end{tabular}
\end{table}

\begin{figure}[H]\centering
\includegraphics[width=0.95\linewidth]{fig_2_3_comparison.png}
\caption{Pair-wise dynamic correlations under EDF, GARCH-N, and GJR-N margins.}
\label{fig:23}
\end{figure}

\paragraph{Reading the comparison.}
The three series share the same pattern, correlations rise
during shocks as the COVID-19 shock (March 2020). Quantitatively, EDF
produces the smoothest path, while GARCH/GJR capture sharper short-run drops in
dependence. The $\hat a+\hat b$ estimate falls from $0.9999$ (EDF) to $\approx
0.96$ once the marginal variance is filtered.

% ---------------------------------------------------------------
\subsection{2.4 -- DCC under a Student-$t$ copula (EDF margins)}
% ---------------------------------------------------------------

We now jointly estimate $(a,b,\nu)$ in \eqref{eq:dcc-LL-t}. The estimates are
$\hat a=0.039216$, $\hat b=0.937845$, $\hat\nu=4.3936$, $\ln L=2093.03$. The
LL gain over the Gaussian-EDF specification is $2093.03-1954.19=138.84$,
overwhelming evidence that the joint tail of the three crypto returns is
heavier than Gaussian and well captured by a $t$-copula with around four
degrees of freedom. Figure~\ref{fig:24} reports the dynamic correlations.

\begin{figure}[H]\centering
\includegraphics[width=0.85\linewidth]{fig_2_4_dcc_t.png}
\caption{Dynamic correlations under DCC-$t$ copula with EDF margins
($\hat\nu=4.39$).}\label{fig:24}
\end{figure}

\paragraph{Intuition for $\hat\nu\approx 4.4$.}
A finite $\nu$ is the device by which the Student-$t$ copula introduces
\emph{tail dependence} (absent in the Gaussian copula). With three crypto
assets sharing common shocks during stress periods, $\hat\nu$ is pushed
toward small values; values in the 4--6 range are typical of equity
indices, and the crypto evidence here lies just at the lower end of that range.
